
\section[Apéndice II: finitud para las imágenes directas de los stacks]
{Apéndice II: teorema de finitud para las imágenes directas de los stacks}
\label{XIII.6}

\begin{proposition}
\label{XIII.6.1}
Sean $S$ un esquema localmente noetheriano y $f\colon X\to S$ un
morfismo. Si $S'$ es un $S$-esquema, se denota por $X'$ (\resp por
$f'$, etc.) la imagen inversa de~$X$ (\resp de~$f$, etc.) por el
morfismo $S'\to S$. Supongamos que, para todo
\marginpar{436}
esquema $S'$ étale sobre~$S$ y todo haz de conjuntos constructible~$F$
sobre~$X'$, $f'_*F$ es constructible, y que, para todo haz de grupos
constructible~$F$ sobre~$X'$, $\R^1f'_*F$ es constructible. Sea $\Phi$
un stack $1$-constructible sobre~$X$~\eqref{XIII.0}. Entonces $f_*\Phi$
es $1$-constructible.
\end{proposition}

Para todo esquema $S'$ étale sobre~$S$ y todo objeto~$x$
de~$(f_*\Phi)_{S'}$, se tiene un isomorfismo
$$
\SheafAut_{S'}(x)\simeq f'_*\SheafAut_{X'}(x)\quoi,
$$
donde, en el segundo miembro de la igualdad, $x$ se considera como
objeto de~$\Phi_{X'}$. Por tanto, las hipótesis implican que $f_*\Phi$
es constructible. Sea $S\Phi$ el haz de los subgerbes maximales
de~$\Phi$~\cite[III~2.1.7]{XIII.1}. Como $f_*(S\Phi)$ es constructible,
puede aplicársele (SGA~4 IX~2.7), y el hecho de que $f_*\Phi$ sea
$1$-constructible se deduce entonces del lema siguiente.

\setcounter{subsection}{1}
\setcounter{subsubsection}{0}
\begin{sublemme}
\label{XIII.6.1.1}
Sean $S$ un esquema localmente noetheriano, $f\colon X\to S$ un
morfismo y $\Phi$ un stack sobre~$X$. Supongamos dados un haz sobre~$S$,
representable por un $S$-esquema~$T$ étale de tipo finito, un morfismo
sobreyectivo
$$
a\colon T\to f_*(S\Phi)
$$
y un objeto~$p$ de la fibra $\Phi_{X_T}$, donde
$X_T=X\times_ST$, que define en
$f_*(S\Phi)(T)=\allowbreak S\Phi(X_T)$ un elemento igual a la imagen~$q$
por~$a$ de la sección identidad de~$T(T)$. Sea
$f_T\colon X_T\to T$ el morfismo canónico, y supongamos que el haz
$\R^1f_{T*}(\SheafAut_{X_T}(p))$ es constructible; entonces también
$S(f_*\Phi)$ es constructible.
\end{sublemme}

El morfismo canónico $f^*f_*\Phi\to\Phi$ da un morfismo
$$
S(f^*f_*\Phi)\simeq f^*(S(f_*\Phi))\to S\Phi,
$$
y por tanto un morfismo canónico
$$
\varphi\colon S(f_*\Phi)\to f_*(S\Phi).
$$
Sean
\marginpar{437}
$F=S(f_*\Phi)$ y $G$ la imagen de~$F$ por~$\varphi$; según
(SGA~4 IX~2.9), $G$ es un haz constructible.

Basta demostrar que, para todo punto~$s$ de~$S$, existe un abierto no
vacío~$U$
\ifthenelse{\boolean{orig}}
{de~$s$}
{de~$\overline{\{s\}}$}
tal que $F|U$ sea localmente constante y constructible.
Sean $s\in S$, $\overline{s}$ un punto geométrico situado sobre~$s$, y
$\overline{q}_1,\dots,\overline{q}_n$ los elementos de~$G_{\overline{s}}$.
Por definición de~$T$, existen $S$-morfismos
\ifthenelse{\boolean{orig}}
{$h_i\colon\overline{s}\to S'$}
{$h_i\colon\overline{s}\to T$}
tales que $\overline{q}_i=h_i^*(q)$. Sean $S'$ el producto fibrado
sobre~$S$ de $n$ esquemas isomorfos a~$T$, $\overline{s}\to S'$ el
producto fibrado de los $h_i$, $X'=X\times_SS'$, y $q_i$ (\resp $p_i$)
la imagen inversa de~$q$ (\resp de~$p$) por la $i$-ésima proyección
de~$S'$ sobre~$T$. Si $\Psi_i$ es el subgerbe maximal de~$\Phi|X'$
engendrado por~$p_i$, el haz
$F_i=\R^1f_*'(\SheafAut_{X'}(p_i))$ no es otro que el haz
$S(f_*'\Psi_i)$ de los subgerbes maximales de~$f_*'\Psi_i$. En
particular, la inyección canónica $\Psi_i\to\Phi|X'$ da un morfismo
$$
\alpha_i\colon F_i\to F|S'.
$$

Vamos a demostrar que $\alpha_i$ es una biyección de~$F_i$ sobre la
imagen inversa de~$q_i$ en~$F|S'$. Para todo esquema $S''$ étale
sobre~$S'$, toda sección~$y$ de~$F_i(S'')$ tiene por imagen
$q_i|S''$ en~$F(S'')$, pues, localmente para la topología étale
sobre~$S''$, $y$ está definida por un objeto~$x$ de~$\Phi_{X''}$
isomorfo a~$p_i|X''$. Recíprocamente, si $y\in F(S'')$ tiene por imagen
$q_i|S''$ en~$F(S'')$, localmente para la topología étale sobre~$S''$,
$y$ está definida por un objeto~$x$ de~$\Phi_{X''}$ isomorfo a~$p_i$;
por consiguiente, $x$ es un objeto de~$\Psi_{iX''}$ y
$y\in F_i(S'')$.

La demostración concluye utilizando~\Ref{XIII.6.1.2}. En efecto, puede
encontrarse un entorno abierto~$U'$ de~$s$ tal que
$q_1|U',\dots,q_n|U'$ sean secciones de~$G(U')$ y engendren este haz.
Como $F_i|U'$ y $G|U'$ son constructibles, también $F|U'$ lo es
según~\Ref{XIII.6.1.2}; tras sustituir, si es preciso, $U$ por un
abierto más pequeño, $F|U$ es localmente constante, lo que concluye la
demostración.
